%% ODER: format ==         = "\mathrel{==}"
%% ODER: format /=         = "\neq "
%
%
\makeatletter
\@ifundefined{lhs2tex.lhs2tex.sty.read}%
  {\@namedef{lhs2tex.lhs2tex.sty.read}{}%
   \newcommand\SkipToFmtEnd{}%
   \newcommand\EndFmtInput{}%
   \long\def\SkipToFmtEnd#1\EndFmtInput{}%
  }\SkipToFmtEnd

\newcommand\ReadOnlyOnce[1]{\@ifundefined{#1}{\@namedef{#1}{}}\SkipToFmtEnd}
\usepackage{amstext}
\usepackage{amssymb}
\usepackage{stmaryrd}
\DeclareFontFamily{OT1}{cmtex}{}
\DeclareFontShape{OT1}{cmtex}{m}{n}
  {<5><6><7><8>cmtex8
   <9>cmtex9
   <10><10.95><12><14.4><17.28><20.74><24.88>cmtex10}{}
\DeclareFontShape{OT1}{cmtex}{m}{it}
  {<-> ssub * cmtt/m/it}{}
\newcommand{\texfamily}{\fontfamily{cmtex}\selectfont}
\DeclareFontShape{OT1}{cmtt}{bx}{n}
  {<5><6><7><8>cmtt8
   <9>cmbtt9
   <10><10.95><12><14.4><17.28><20.74><24.88>cmbtt10}{}
\DeclareFontShape{OT1}{cmtex}{bx}{n}
  {<-> ssub * cmtt/bx/n}{}
\newcommand{\tex}[1]{\text{\texfamily#1}}	% NEU

\newcommand{\Sp}{\hskip.33334em\relax}


\newcommand{\Conid}[1]{\mathit{#1}}
\newcommand{\Varid}[1]{\mathit{#1}}
\newcommand{\anonymous}{\kern0.06em \vbox{\hrule\@width.5em}}
\newcommand{\plus}{\mathbin{+\!\!\!+}}
\newcommand{\bind}{\mathbin{>\!\!\!>\mkern-6.7mu=}}
\newcommand{\rbind}{\mathbin{=\mkern-6.7mu<\!\!\!<}}% suggested by Neil Mitchell
\newcommand{\sequ}{\mathbin{>\!\!\!>}}
\renewcommand{\leq}{\leqslant}
\renewcommand{\geq}{\geqslant}
\usepackage{polytable}

%mathindent has to be defined
\@ifundefined{mathindent}%
  {\newdimen\mathindent\mathindent\leftmargini}%
  {}%

\def\resethooks{%
  \global\let\SaveRestoreHook\empty
  \global\let\ColumnHook\empty}
\newcommand*{\savecolumns}[1][default]%
  {\g@addto@macro\SaveRestoreHook{\savecolumns[#1]}}
\newcommand*{\restorecolumns}[1][default]%
  {\g@addto@macro\SaveRestoreHook{\restorecolumns[#1]}}
\newcommand*{\aligncolumn}[2]%
  {\g@addto@macro\ColumnHook{\column{#1}{#2}}}

\resethooks

\newcommand{\onelinecommentchars}{\quad-{}- }
\newcommand{\commentbeginchars}{\enskip\{-}
\newcommand{\commentendchars}{-\}\enskip}

\newcommand{\visiblecomments}{%
  \let\onelinecomment=\onelinecommentchars
  \let\commentbegin=\commentbeginchars
  \let\commentend=\commentendchars}

\newcommand{\invisiblecomments}{%
  \let\onelinecomment=\empty
  \let\commentbegin=\empty
  \let\commentend=\empty}

\visiblecomments

\newlength{\blanklineskip}
\setlength{\blanklineskip}{0.66084ex}

\newcommand{\hsindent}[1]{\quad}% default is fixed indentation
\let\hspre\empty
\let\hspost\empty
\newcommand{\NB}{\textbf{NB}}
\newcommand{\Todo}[1]{$\langle$\textbf{To do:}~#1$\rangle$}

\EndFmtInput
\makeatother
%
%
%
%
%
%
% This package provides two environments suitable to take the place
% of hscode, called "plainhscode" and "arrayhscode". 
%
% The plain environment surrounds each code block by vertical space,
% and it uses \abovedisplayskip and \belowdisplayskip to get spacing
% similar to formulas. Note that if these dimensions are changed,
% the spacing around displayed math formulas changes as well.
% All code is indented using \leftskip.
%
% Changed 19.08.2004 to reflect changes in colorcode. Should work with
% CodeGroup.sty.
%
\ReadOnlyOnce{polycode.fmt}%
\makeatletter

\newcommand{\hsnewpar}[1]%
  {{\parskip=0pt\parindent=0pt\par\vskip #1\noindent}}

% can be used, for instance, to redefine the code size, by setting the
% command to \small or something alike
\newcommand{\hscodestyle}{}

% The command \sethscode can be used to switch the code formatting
% behaviour by mapping the hscode environment in the subst directive
% to a new LaTeX environment.

\newcommand{\sethscode}[1]%
  {\expandafter\let\expandafter\hscode\csname #1\endcsname
   \expandafter\let\expandafter\endhscode\csname end#1\endcsname}

% "compatibility" mode restores the non-polycode.fmt layout.

\newenvironment{compathscode}%
  {\par\noindent
   \advance\leftskip\mathindent
   \hscodestyle
   \let\\=\@normalcr
   \let\hspre\(\let\hspost\)%
   \pboxed}%
  {\endpboxed\)%
   \par\noindent
   \ignorespacesafterend}

\newcommand{\compaths}{\sethscode{compathscode}}

% "plain" mode is the proposed default.
% It should now work with \centering.
% This required some changes. The old version
% is still available for reference as oldplainhscode.

\newenvironment{plainhscode}%
  {\hsnewpar\abovedisplayskip
   \advance\leftskip\mathindent
   \hscodestyle
   \let\hspre\(\let\hspost\)%
   \pboxed}%
  {\endpboxed%
   \hsnewpar\belowdisplayskip
   \ignorespacesafterend}

\newenvironment{oldplainhscode}%
  {\hsnewpar\abovedisplayskip
   \advance\leftskip\mathindent
   \hscodestyle
   \let\\=\@normalcr
   \(\pboxed}%
  {\endpboxed\)%
   \hsnewpar\belowdisplayskip
   \ignorespacesafterend}

% Here, we make plainhscode the default environment.

\newcommand{\plainhs}{\sethscode{plainhscode}}
\newcommand{\oldplainhs}{\sethscode{oldplainhscode}}
\plainhs

% The arrayhscode is like plain, but makes use of polytable's
% parray environment which disallows page breaks in code blocks.

\newenvironment{arrayhscode}%
  {\hsnewpar\abovedisplayskip
   \advance\leftskip\mathindent
   \hscodestyle
   \let\\=\@normalcr
   \(\parray}%
  {\endparray\)%
   \hsnewpar\belowdisplayskip
   \ignorespacesafterend}

\newcommand{\arrayhs}{\sethscode{arrayhscode}}

% The mathhscode environment also makes use of polytable's parray 
% environment. It is supposed to be used only inside math mode 
% (I used it to typeset the type rules in my thesis).

\newenvironment{mathhscode}%
  {\parray}{\endparray}

\newcommand{\mathhs}{\sethscode{mathhscode}}

% texths is similar to mathhs, but works in text mode.

\newenvironment{texthscode}%
  {\(\parray}{\endparray\)}

\newcommand{\texths}{\sethscode{texthscode}}

% The framed environment places code in a framed box.

\def\codeframewidth{\arrayrulewidth}
\RequirePackage{calc}

\newenvironment{framedhscode}%
  {\parskip=\abovedisplayskip\par\noindent
   \hscodestyle
   \arrayrulewidth=\codeframewidth
   \tabular{@{}|p{\linewidth-2\arraycolsep-2\arrayrulewidth-2pt}|@{}}%
   \hline\framedhslinecorrect\\{-1.5ex}%
   \let\endoflinesave=\\
   \let\\=\@normalcr
   \(\pboxed}%
  {\endpboxed\)%
   \framedhslinecorrect\endoflinesave{.5ex}\hline
   \endtabular
   \parskip=\belowdisplayskip\par\noindent
   \ignorespacesafterend}

\newcommand{\framedhslinecorrect}[2]%
  {#1[#2]}

\newcommand{\framedhs}{\sethscode{framedhscode}}

% The inlinehscode environment is an experimental environment
% that can be used to typeset displayed code inline.

\newenvironment{inlinehscode}%
  {\(\def\column##1##2{}%
   \let\>\undefined\let\<\undefined\let\\\undefined
   \newcommand\>[1][]{}\newcommand\<[1][]{}\newcommand\\[1][]{}%
   \def\fromto##1##2##3{##3}%
   \def\nextline{}}{\) }%

\newcommand{\inlinehs}{\sethscode{inlinehscode}}

% The joincode environment is a separate environment that
% can be used to surround and thereby connect multiple code
% blocks.

\newenvironment{joincode}%
  {\let\orighscode=\hscode
   \let\origendhscode=\endhscode
   \def\endhscode{\def\hscode{\endgroup\def\@currenvir{hscode}\\}\begingroup}
   %\let\SaveRestoreHook=\empty
   %\let\ColumnHook=\empty
   %\let\resethooks=\empty
   \orighscode\def\hscode{\endgroup\def\@currenvir{hscode}}}%
  {\origendhscode
   \global\let\hscode=\orighscode
   \global\let\endhscode=\origendhscode}%

\makeatother
\EndFmtInput
%

% format { = "\lbag"
% format } = "\rbag"


% TODO: fix spacing for \lbag and \rbag





















































%formap phiX = "\Phi"














\section{Bidirectional programming on lists}

\ignore{

\begin{hscode}\SaveRestoreHook
\column{B}{@{}>{\hspre}l<{\hspost}@{}}%
\column{E}{@{}>{\hspre}l<{\hspost}@{}}%
\>[B]{}\mbox{\enskip\{-\# LANGUAGE FlexibleContexts, TemplateHaskell, TypeFamilies  \#-\}\enskip}{}\<[E]%
\\[\blanklineskip]%
\>[B]{}\mathbf{module}\;\Conid{PList}\;\mathbf{where}{}\<[E]%
\\
\>[B]{}\mathbf{import}\;\Conid{\Conid{Generics}.BiGUL}{}\<[E]%
\\
\>[B]{}\mathbf{import}\;\Conid{\Conid{Generics}.\Conid{BiGUL}.Interpreter}{}\<[E]%
\\
\>[B]{}\mathbf{import}\;\Conid{\Conid{Generics}.\Conid{BiGUL}.TH}{}\<[E]%
\\
\>[B]{}\mathbf{import}\;\Conid{\Conid{Generics}.\Conid{BiGUL}.Lib}{}\<[E]%
\\
\>[B]{}\mathbf{import}\;\Conid{\Conid{Data}.List}{}\<[E]%
\\
\>[B]{}\mathbf{import}\;\Conid{\Conid{Data}.Maybe}{}\<[E]%
\\
\>[B]{}\mathbf{import}\;\Conid{\Conid{Control}.\Conid{Monad}.Except}{}\<[E]%
\\
\>[B]{}\mathbf{import}\;\Conid{\Conid{GHC}.Generics}{}\<[E]%
\\
\>[B]{}\mathbf{import}\;\Conid{PBasic}{}\<[E]%
\ColumnHook
\end{hscode}\resethooks

}

In this section, we demosntrate that
many list functions can be bidirectionalized
using BiGUL. To show the correspondence with the original list functions,
we prefix
the original forward function names with "lens". Note that in our context,
the original forward
functions can be automatically derived from the new putback transformations
by calling \ensuremath{\Varid{get}}.

We shall focus on bidirectionalize \ensuremath{\Varid{foldr}}, a simple but useful higher order function on lists:\begin{hscode}\SaveRestoreHook
\column{B}{@{}>{\hspre}l<{\hspost}@{}}%
\column{3}{@{}>{\hspre}l<{\hspost}@{}}%
\column{21}{@{}>{\hspre}l<{\hspost}@{}}%
\column{E}{@{}>{\hspre}l<{\hspost}@{}}%
\>[3]{}\Varid{foldr}\;\Varid{f}\;\Varid{e}\;[\mskip1.5mu \mskip1.5mu]{}\<[21]%
\>[21]{}\mathrel{=}\Varid{e}{}\<[E]%
\\
\>[3]{}\Varid{foldr}\;\Varid{f}\;\Varid{e}\;(\Varid{x}\mathbin{:}\Varid{xs}){}\<[21]%
\>[21]{}\mathrel{=}\Varid{f}\;\Varid{x}\;(\Varid{foldr}\;\Varid{f}\;\Varid{e}\;\Varid{xs}){}\<[E]%
\ColumnHook
\end{hscode}\resethooks
where many interesting functions can be defined in terms of \ensuremath{\Varid{foldr}}:\begin{hscode}\SaveRestoreHook
\column{B}{@{}>{\hspre}l<{\hspost}@{}}%
\column{3}{@{}>{\hspre}l<{\hspost}@{}}%
\column{E}{@{}>{\hspre}l<{\hspost}@{}}%
\>[3]{}\Varid{sum}\mathrel{=}\Varid{foldr}\;(\mathbin{+})\;\mathrm{0}{}\<[E]%
\\
\>[3]{}\Varid{map}\;\Varid{f}\mathrel{=}\Varid{foldr}\;(\lambda \Varid{a}\;\Varid{r}\to \Varid{f}\;\Varid{a}\mathbin{:}\Varid{r})\;[\mskip1.5mu \mskip1.5mu]{}\<[E]%
\\
\>[3]{}\Varid{filter}\;\Varid{p}\mathrel{=}\Varid{foldr}\;(\lambda \Varid{a}\;\Varid{r}\to \mathbf{if}\;\Varid{p}\;\Varid{a}\;\mathbf{then}\;\Varid{a}\mathbin{:}\Varid{r}\;\mathbf{else}\;\Varid{r})\;[\mskip1.5mu \mskip1.5mu]{}\<[E]%
\\
\>[3]{}\Varid{reverse}\;\Varid{p}\mathrel{=}\Varid{foldr}\;(\lambda \Varid{a}\;\Varid{r}\to \Varid{r}\plus [\mskip1.5mu \Varid{a}\mskip1.5mu])\;[\mskip1.5mu \mskip1.5mu]{}\<[E]%
\\
\>[3]{}\Varid{sort}\mathrel{=}\Varid{foldr}\;\Varid{insert}\;[\mskip1.5mu \mskip1.5mu]{}\<[E]%
\ColumnHook
\end{hscode}\resethooks
We start by developing a putback version for \ensuremath{\Varid{foldr}}.
\begin{hscode}\SaveRestoreHook
\column{B}{@{}>{\hspre}l<{\hspost}@{}}%
\column{3}{@{}>{\hspre}l<{\hspost}@{}}%
\column{9}{@{}>{\hspre}c<{\hspost}@{}}%
\column{9E}{@{}l@{}}%
\column{12}{@{}>{\hspre}l<{\hspost}@{}}%
\column{14}{@{}>{\hspre}l<{\hspost}@{}}%
\column{15}{@{}>{\hspre}l<{\hspost}@{}}%
\column{19}{@{}>{\hspre}l<{\hspost}@{}}%
\column{21}{@{}>{\hspre}l<{\hspost}@{}}%
\column{25}{@{}>{\hspre}l<{\hspost}@{}}%
\column{60}{@{}>{\hspre}l<{\hspost}@{}}%
\column{E}{@{}>{\hspre}l<{\hspost}@{}}%
\>[B]{}\Varid{lensFoldr}\mathbin{::}{}\<[15]%
\>[15]{}(\Conid{Show}\;\Varid{a},\Conid{Show}\;\Varid{b})\Rightarrow {}\<[E]%
\\
\>[15]{}\Conid{BiGUL}\;(\Varid{a},\Varid{b})\;\Varid{b}\to (\Varid{b}\to \Conid{Bool})\to \Conid{BiGUL}\;([\mskip1.5mu \Varid{a}\mskip1.5mu],\Varid{b})\;\Varid{b}{}\<[E]%
\\
\>[B]{}\Varid{lensFoldr}\;\Varid{bx}\;\Varid{pv}\mathrel{=}{}\<[E]%
\\
\>[B]{}\hsindent{3}{}\<[3]%
\>[3]{}\Conid{Case}\;{}\<[9]%
\>[9]{}[\mskip1.5mu {}\<[9E]%
\>[12]{}\mathbin{\$}(\Varid{adaptive}\;[\mskip1.5mu \mid \lambda (\Varid{x},\Varid{y})\;\Varid{v}\to \Varid{pv}\;\Varid{v}\mathrel{\wedge}\Varid{length}\;\Varid{x}\not\equiv \mathrm{0}\mid \mskip1.5mu]){}\<[E]%
\\
\>[12]{}\hsindent{2}{}\<[14]%
\>[14]{}\Longrightarrow\lambda (\Varid{x},\Varid{y})\;\Varid{v}\to ([\mskip1.5mu \mskip1.5mu],\Varid{y}){}\<[E]%
\\
\>[9]{},{}\<[9E]%
\>[12]{}\mathbin{\$}(\Varid{normal}\;[\mskip1.5mu \mid \lambda (\Varid{xs},\anonymous )\;\Varid{v}\to \Varid{null}\;\Varid{xs}\mid \mskip1.5mu]\;[\mskip1.5mu \mid \lambda (\Varid{xs},\anonymous )\to \Varid{null}\;\Varid{xs}\mid \mskip1.5mu]){}\<[E]%
\\
\>[12]{}\hsindent{2}{}\<[14]%
\>[14]{}\Longrightarrow{}\<[19]%
\>[19]{}\mathbin{\$}(\Varid{rearrV}\;[\mskip1.5mu \mid \lambda \Varid{v}\to ((),\Varid{v})\mid \mskip1.5mu])\mathbin{\$}{}\<[E]%
\\
\>[19]{}\hsindent{6}{}\<[25]%
\>[25]{}\mathbin{\$}(\Varid{update}\;[\mskip1.5mu \Varid{p}\mid (\anonymous ,\Varid{v})\mid \mskip1.5mu]\;[\mskip1.5mu \Varid{p}\mid ((),\Varid{v})\mid \mskip1.5mu]\;[\mskip1.5mu \Varid{d}\mid \Varid{v}\mathrel{=}\Conid{Replace}\mid \mskip1.5mu]){}\<[E]%
\\
\>[9]{},{}\<[9E]%
\>[12]{}\mathbin{\$}(\Varid{normalSV}\;[\mskip1.5mu \Varid{p}\mid \anonymous \mid \mskip1.5mu]\;[\mskip1.5mu \Varid{p}\mid \anonymous \mid \mskip1.5mu]\;[\mskip1.5mu \mid \lambda (\Varid{xs},\anonymous )\to not\;(\Varid{null}\;\Varid{xs})\mid \mskip1.5mu]){}\<[E]%
\\
\>[12]{}\hsindent{2}{}\<[14]%
\>[14]{}\Longrightarrow{}\<[19]%
\>[19]{}\mathbin{\$}(\Varid{rearrS}\;[\mskip1.5mu \mid \lambda ((\Varid{x}\mathbin{:}\Varid{xs}),\Varid{e})\to (\Varid{x},(\Varid{xs},\Varid{e})){}\<[60]%
\>[60]{}\mid \mskip1.5mu])\;{}\<[E]%
\\
\>[19]{}\hsindent{2}{}\<[21]%
\>[21]{}(\Conid{Replace}\mathbin{`\Conid{Prod}`}\Varid{lensFoldr}\;\Varid{bx}\;\Varid{pv})\mathbin{`\Conid{Compose}`}\Varid{bx}{}\<[E]%
\\
\>[9]{}\mskip1.5mu]{}\<[9E]%
\ColumnHook
\end{hscode}\resethooks
The above \ensuremath{\Varid{lensFoldr}\;\Varid{bx}\;\Varid{pv}} updates source \ensuremath{(\Varid{xs},\Varid{y})} with view \ensuremath{\Varid{v}}.
If \ensuremath{\Varid{v}} satisfies \ensuremath{\Varid{pv}}, it will try to stop by adapting \ensuremath{\Varid{xs}} to \ensuremath{[\mskip1.5mu \mskip1.5mu]}.
Or it will recursively apply \ensuremath{\Varid{bx}} rightwards over
the elements of \ensuremath{\Varid{xs}}: if \ensuremath{\Varid{xs}} is empty, it will just embed \ensuremath{\Varid{v}} to \ensuremath{\Varid{y}};
otherwise it will rearrange the source for the recursive call.
\ignore{
lensFilter :: (a->Bool) -> BiGUL [a] [a]
lensInsert :: Ord a => BiGUL (a,[a]{-sorted-}) [a]{-sorted-}
lensSort :: Ord a => BiGUL [a] [a]
}

With \ensuremath{\Varid{lensFoldr}}, we can redefine many list functions from the putback point
of view. As the first example, \ensuremath{\Varid{map}} can be defined in terms of \ensuremath{\Varid{lensFoldr}}.

\begin{hscode}\SaveRestoreHook
\column{B}{@{}>{\hspre}l<{\hspost}@{}}%
\column{4}{@{}>{\hspre}l<{\hspost}@{}}%
\column{11}{@{}>{\hspre}l<{\hspost}@{}}%
\column{15}{@{}>{\hspre}l<{\hspost}@{}}%
\column{18}{@{}>{\hspre}l<{\hspost}@{}}%
\column{21}{@{}>{\hspre}l<{\hspost}@{}}%
\column{E}{@{}>{\hspre}l<{\hspost}@{}}%
\>[B]{}\Varid{lensMap}\mathbin{::}(\Conid{Show}\;\Varid{a},\Conid{Show}\;\Varid{b})\Rightarrow \Conid{BiGUL}\;\Varid{a}\;\Varid{b}\to \Conid{BiGUL}\;([\mskip1.5mu \Varid{a}\mskip1.5mu],[\mskip1.5mu \Varid{b}\mskip1.5mu])\;[\mskip1.5mu \Varid{b}\mskip1.5mu]{}\<[E]%
\\
\>[B]{}\Varid{lensMap}\;\Varid{bx}\mathrel{=}{}\<[15]%
\>[15]{}\Varid{lensFoldr}\;\Varid{bx'}\;\Varid{null}{}\<[E]%
\\
\>[B]{}\hsindent{4}{}\<[4]%
\>[4]{}\mathbf{where}\;{}\<[11]%
\>[11]{}\Varid{bx'}\mathrel{=}{}\<[18]%
\>[18]{}\mathbin{\$}(\Varid{rearrV}\;[\mskip1.5mu \mid \lambda (\Varid{v}\mathbin{:}\Varid{vs})\to (\Varid{v},\Varid{vs})\mid \mskip1.5mu])\mathbin{\$}{}\<[E]%
\\
\>[18]{}\hsindent{3}{}\<[21]%
\>[21]{}\Varid{bx}\mathbin{`\Conid{Prod}`}\Conid{Replace}{}\<[E]%
\ColumnHook
\end{hscode}\resethooks
\begin{tabbing}\tt
~\char42{}PList\char62{}~put~\char40{}lensMap~dec1\char41{}~\char40{}\char91{}0\char46{}\char46{}10\char93{}\char44{}\char91{}\char93{}\char41{}~\char91{}100\char46{}\char46{}110\char93{}\\
\tt ~Just~\char40{}\char91{}99\char44{}100\char44{}101\char44{}102\char44{}103\char44{}104\char44{}105\char44{}106\char44{}107\char44{}108\char44{}109\char93{}\char44{}\char91{}\char93{}\char41{}\\
\tt ~\char42{}PList\char62{}~get~\char40{}lensMap~dec1\char41{}~\char40{}\char91{}1\char46{}\char46{}10\char93{}\char44{}\char91{}\char93{}\char41{}\\
\tt ~Just~\char91{}2\char44{}3\char44{}4\char44{}5\char44{}6\char44{}7\char44{}8\char44{}9\char44{}10\char44{}11\char93{}
\end{tabbing}
Here, for testing, we embed into our framework
the bijective functions for increasing and decreasing an integer by \ensuremath{\mathrm{1}}.
\begin{hscode}\SaveRestoreHook
\column{B}{@{}>{\hspre}l<{\hspost}@{}}%
\column{3}{@{}>{\hspre}l<{\hspost}@{}}%
\column{10}{@{}>{\hspre}l<{\hspost}@{}}%
\column{E}{@{}>{\hspre}l<{\hspost}@{}}%
\>[B]{}\Varid{dec1}\mathbin{::}\Conid{BiGUL}\;\mathit{Int}\;\mathit{Int}{}\<[E]%
\\
\>[B]{}\Varid{dec1}\mathrel{=}\Varid{emb}\;\Varid{g}\;\Varid{p}{}\<[E]%
\\
\>[B]{}\hsindent{3}{}\<[3]%
\>[3]{}\mathbf{where}\;{}\<[10]%
\>[10]{}\Varid{g}\;\Varid{s}\mathrel{=}\Varid{s}\mathbin{+}\mathrm{1}{}\<[E]%
\\
\>[10]{}\Varid{p}\;\Varid{s}\;\Varid{v}\mathrel{=}\Varid{v}\mathbin{-}\mathrm{1}{}\<[E]%
\ColumnHook
\end{hscode}\resethooks
 
The second example is to bidirectinalize \ensuremath{\Varid{reverse}}, which is to reverse the
elements of a list.
\begin{hscode}\SaveRestoreHook
\column{B}{@{}>{\hspre}l<{\hspost}@{}}%
\column{16}{@{}>{\hspre}l<{\hspost}@{}}%
\column{22}{@{}>{\hspre}c<{\hspost}@{}}%
\column{22E}{@{}l@{}}%
\column{25}{@{}>{\hspre}l<{\hspost}@{}}%
\column{27}{@{}>{\hspre}l<{\hspost}@{}}%
\column{33}{@{}>{\hspre}l<{\hspost}@{}}%
\column{E}{@{}>{\hspre}l<{\hspost}@{}}%
\>[B]{}\Varid{lensReverse}\mathbin{::}\Conid{Show}\;\Varid{a}\Rightarrow \Conid{BiGUL}\;[\mskip1.5mu \Varid{a}\mskip1.5mu]\;[\mskip1.5mu \Varid{a}\mskip1.5mu]{}\<[E]%
\\
\>[B]{}\Varid{lensReverse}\mathrel{=}{}\<[16]%
\>[16]{}\Conid{Case}\;{}\<[22]%
\>[22]{}[\mskip1.5mu {}\<[22E]%
\\
\>[22]{}\hsindent{3}{}\<[25]%
\>[25]{}\mathbin{\$}(\Varid{adaptive}\;[\mskip1.5mu \mid \lambda \Varid{s}\;\Varid{v}\to \Varid{length}\;\Varid{s}\mathbin{<}\Varid{length}\;\Varid{v}\mid \mskip1.5mu]){}\<[E]%
\\
\>[25]{}\hsindent{2}{}\<[27]%
\>[27]{}\Longrightarrow\lambda \Varid{s}\;\Varid{v}\to \Varid{v}{}\<[E]%
\\
\>[22]{},{}\<[22E]%
\>[25]{}\mathbin{\$}(\Varid{normalSV}\;[\mskip1.5mu \Varid{p}\mid \anonymous \mid \mskip1.5mu]\;[\mskip1.5mu \Varid{p}\mid \anonymous \mid \mskip1.5mu]\;[\mskip1.5mu \mid \lambda \Varid{s}\to \Conid{True}\mid \mskip1.5mu]){}\<[E]%
\\
\>[25]{}\hsindent{2}{}\<[27]%
\>[27]{}\Longrightarrow\mathbin{\$}(\Varid{rearrS}\;[\mskip1.5mu \mid \lambda \Varid{s}\to (\Varid{s},[\mskip1.5mu \mskip1.5mu])\mid \mskip1.5mu])\mathbin{\$}{}\<[E]%
\\
\>[27]{}\hsindent{6}{}\<[33]%
\>[33]{}\Varid{lensFoldr}\;(\Varid{lensSwap}\mathbin{`\Conid{Compose}`}\Varid{lensSnoc})\;\Varid{null}{}\<[E]%
\\
\>[22]{}\mskip1.5mu]{}\<[22E]%
\ColumnHook
\end{hscode}\resethooks
When the source is shorter than the view, we add more elements to the source
by filling in it with the elements of the view (Note that this is one option;
we could do other ways by, say, duplicating elements in the source). When the source is long enough, we decompose the view into a snoc-list using \ensuremath{\Varid{lensSnoc}}, and use it
to update the cons-list of the source.
\begin{hscode}\SaveRestoreHook
\column{B}{@{}>{\hspre}l<{\hspost}@{}}%
\column{18}{@{}>{\hspre}c<{\hspost}@{}}%
\column{18E}{@{}l@{}}%
\column{21}{@{}>{\hspre}l<{\hspost}@{}}%
\column{24}{@{}>{\hspre}l<{\hspost}@{}}%
\column{31}{@{}>{\hspre}l<{\hspost}@{}}%
\column{35}{@{}>{\hspre}l<{\hspost}@{}}%
\column{E}{@{}>{\hspre}l<{\hspost}@{}}%
\>[B]{}\Varid{lensSnoc}\mathbin{::}\Conid{Show}\;\Varid{a}\Rightarrow \Conid{BiGUL}\;([\mskip1.5mu \Varid{a}\mskip1.5mu],\Varid{a})\;[\mskip1.5mu \Varid{a}\mskip1.5mu]{}\<[E]%
\\
\>[B]{}\Varid{lensSnoc}\mathrel{=}\Conid{Case}\;{}\<[18]%
\>[18]{}[\mskip1.5mu {}\<[18E]%
\>[21]{}\mathbin{\$}(\Varid{normal}\;[\mskip1.5mu \mid \lambda \Varid{s}\;\Varid{v}\to \Varid{length}\;\Varid{v}\equiv \mathrm{1}\mid \mskip1.5mu]\;[\mskip1.5mu \mid \lambda (\Varid{s},\anonymous )\to \Varid{null}\;\Varid{s}\mid \mskip1.5mu]){}\<[E]%
\\
\>[21]{}\hsindent{3}{}\<[24]%
\>[24]{}\Longrightarrow\mathbin{\$}(\Varid{rearrV}\;[\mskip1.5mu \mid \lambda [\mskip1.5mu \Varid{v}\mskip1.5mu]\to ([\mskip1.5mu \mskip1.5mu],\Varid{v})\mid \mskip1.5mu])\;\Conid{Replace}{}\<[E]%
\\
\>[18]{},{}\<[18E]%
\>[21]{}\mathbin{\$}(\Varid{normal}\;[\mskip1.5mu \mid \lambda (\Varid{s},\anonymous )\;\Varid{v}\to \Varid{length}\;\Varid{s}\mathbin{>}\mathrm{0}\mid \mskip1.5mu]\;[\mskip1.5mu \mid \lambda (\Varid{s},\anonymous )\to \Varid{length}\;\Varid{s}\mathbin{>}\mathrm{0}\mid \mskip1.5mu]){}\<[E]%
\\
\>[21]{}\hsindent{3}{}\<[24]%
\>[24]{}\Longrightarrow\mathbin{\$}(\Varid{rearrS}\;[\mskip1.5mu \mid \lambda (\Varid{y}\mathbin{:}\Varid{ys},\Varid{x})\to (\Varid{y},(\Varid{ys},\Varid{x}))\mid \mskip1.5mu])\mathbin{\$}{}\<[E]%
\\
\>[24]{}\hsindent{7}{}\<[31]%
\>[31]{}\mathbin{\$}(\Varid{rearrV}\;[\mskip1.5mu \mid \lambda (\Varid{v}\mathbin{:}\Varid{vs})\to (\Varid{v},\Varid{vs})\mid \mskip1.5mu])\mathbin{\$}{}\<[E]%
\\
\>[31]{}\hsindent{4}{}\<[35]%
\>[35]{}\Conid{Replace}\mathbin{`\Conid{Prod}`}\Varid{lensSnoc}{}\<[E]%
\\
\>[18]{},{}\<[18E]%
\>[21]{}\mathbin{\$}(\Varid{adaptive}\;[\mskip1.5mu \mid \lambda (\Varid{s},\anonymous )\;\Varid{v}\to \Varid{null}\;\Varid{s}\mid \mskip1.5mu]){}\<[E]%
\\
\>[21]{}\hsindent{3}{}\<[24]%
\>[24]{}\Longrightarrow\lambda (\Varid{s},\Varid{x})\;\anonymous \to ([\mskip1.5mu \bot \mskip1.5mu],\Varid{x}){}\<[E]%
\\
\>[18]{}\mskip1.5mu]{}\<[18E]%
\\[\blanklineskip]%
\>[B]{}\Varid{lensSwap}\mathbin{::}(\Conid{Show}\;\Varid{a},\Conid{Show}\;\Varid{b})\Rightarrow \Conid{BiGUL}\;(\Varid{a},\Varid{b})\;(\Varid{b},\Varid{a}){}\<[E]%
\\
\>[B]{}\Varid{lensSwap}\mathrel{=}\mathbin{\$}(\Varid{rearrS}\;[\mskip1.5mu \mid \lambda (\Varid{x},\Varid{y})\to (\Varid{y},\Varid{x})\mid \mskip1.5mu])\;\Conid{Replace}{}\<[E]%
\ColumnHook
\end{hscode}\resethooks

Below are some testing examples.
\begin{tabbing}\tt
~\char42{}PList\char62{}~put~lensSnoc~\char40{}\char91{}2\char44{}3\char44{}4\char93{}\char44{}1\char41{}~\char91{}10\char44{}11\char44{}12\char44{}13\char93{}\\
\tt ~Just~\char40{}\char91{}10\char44{}11\char44{}12\char93{}\char44{}13\char41{}\\
\tt ~\char42{}PList\char62{}~put~lensSnoc~\char40{}\char91{}2\char44{}3\char44{}4\char93{}\char44{}1\char41{}~\char91{}10\char44{}11\char44{}12\char44{}13\char44{}14\char93{}\\
\tt ~Just~\char40{}\char91{}10\char44{}11\char44{}12\char44{}13\char93{}\char44{}14\char41{}\\
\tt ~\char42{}PList\char62{}~put~lensSnoc~\char40{}\char91{}2\char44{}3\char44{}4\char93{}\char44{}1\char41{}~\char91{}10\char44{}11\char93{}\\
\tt ~Just~\char40{}\char91{}10\char93{}\char44{}11\char41{}\\
\tt ~\char42{}PList\char62{}~get~lensSnoc~\char40{}\char91{}1\char46{}\char46{}10\char93{}\char44{}~100\char41{}\\
\tt ~Just~\char91{}1\char44{}2\char44{}3\char44{}4\char44{}5\char44{}6\char44{}7\char44{}8\char44{}9\char44{}10\char44{}100\char93{}\\
\tt ~\\
\tt ~\char42{}PList\char62{}~put~lensReverse~\char91{}1\char46{}\char46{}10\char93{}~\char91{}100\char46{}\char46{}105\char93{}\\
\tt ~Just~\char91{}105\char44{}104\char44{}103\char44{}102\char44{}101\char44{}100\char93{}\\
\tt ~\char42{}PList\char62{}~put~lensReverse~\char91{}1\char46{}\char46{}10\char93{}~\char91{}100\char46{}\char46{}115\char93{}\\
\tt ~Just~\char91{}115\char44{}114\char44{}113\char44{}112\char44{}111\char44{}110\char44{}109\char44{}108\char44{}107\char44{}106\char44{}105\char44{}104\char44{}103\char44{}102\char44{}101\char44{}100\char93{}\\
\tt ~\char42{}PList\char62{}~get~lensReverse~\char91{}1\char46{}\char46{}10\char93{}\\
\tt ~Just~\char91{}10\char44{}9\char44{}8\char44{}7\char44{}6\char44{}5\char44{}4\char44{}3\char44{}2\char44{}1\char93{}
\end{tabbing}

It is worth noting that our definition of \ensuremath{\Varid{lensFoldr}} is just one
candidate putback function for \ensuremath{\Varid{foldr}}, and there are many others.
This reflects the fact
that one \ensuremath{\Varid{foldr}} can have many \ensuremath{\Varid{put}}s, each describing different
updating strategy.

\ignore{
\subsection{Efficiency Issue: \ensuremath{\Varid{lensFoldr}}}

A close look at the definition of \ensuremath{\Varid{lensFoldr}} reveals that it contains many
redundant computations because of the use of \ensuremath{\Conid{Compose}} that calls \ensuremath{\Varid{get}} as many
times as the number of calls of \ensuremath{\Conid{Compose}}. Put it more concretely,\begin{hscode}\SaveRestoreHook
\column{B}{@{}>{\hspre}l<{\hspost}@{}}%
\column{3}{@{}>{\hspre}l<{\hspost}@{}}%
\column{E}{@{}>{\hspre}l<{\hspost}@{}}%
\>[3]{}\Varid{put}\;(\Varid{lensFoldr}\;\Varid{bx}\;(\Varid{const}\;\Conid{True}))\;([\mskip1.5mu \Varid{x1},\Varid{x2},\mathbin{...},\Varid{xn}\mskip1.5mu],\Varid{e})\;\Varid{v}{}\<[E]%
\ColumnHook
\end{hscode}\resethooks
would call\begin{hscode}\SaveRestoreHook
\column{B}{@{}>{\hspre}l<{\hspost}@{}}%
\column{3}{@{}>{\hspre}l<{\hspost}@{}}%
\column{E}{@{}>{\hspre}l<{\hspost}@{}}%
\>[3]{}\Varid{get}\;(\Varid{lensFoldr}\;\Varid{bx}\;(\Varid{const}\;\Conid{True}))\;([\mskip1.5mu \Varid{x2},\mathbin{...},\Varid{xn}\mskip1.5mu],\Varid{e}),{}\<[E]%
\\
\>[3]{}\mathbin{...},{}\<[E]%
\\
\>[3]{}\Varid{get}\;(\Varid{lensFoldr}\;\Varid{bx}\;(\Varid{const}\;\Conid{True}))\;([\mskip1.5mu \mskip1.5mu],\Varid{e})\mathbin{\circ}{}\<[E]%
\ColumnHook
\end{hscode}\resethooks

\subsection{LensScanr}

\subsection{LensMap}
}
